%% Manju Vallayil — General CV (Data Science / AI Research)
%% Compile with: pdflatex general_cv.tex  (run twice)

\documentclass[10.5pt,a4paper]{article}

% ── Packages ──────────────────────────────────────────────────────────────────
\usepackage[T1]{fontenc}
\usepackage[utf8]{inputenc}
\usepackage[left=1.8cm, right=1.8cm, top=1.5cm, bottom=1.8cm]{geometry}
\usepackage{charter}
\usepackage{microtype}
\usepackage{parskip}
\usepackage{enumitem}
\usepackage{xcolor}
\usepackage{titlesec}
\usepackage{array}
\usepackage{tabularx}
\usepackage[most]{tcolorbox}
\usepackage[
  colorlinks=true,
  urlcolor=accent,
  linkcolor=accent,
  citecolor=accent
]{hyperref}
\usepackage{fontawesome5}

% ── Colours ───────────────────────────────────────────────────────────────────
\definecolor{accent}{HTML}{3E0097}
\definecolor{accentlt}{HTML}{EDE8F7}
\definecolor{mid}{HTML}{444444}
\definecolor{light}{HTML}{999999}

% ── Section headings ──────────────────────────────────────────────────────────
\titleformat{\section}
  {\color{accent}\small\bfseries\scshape}{}{0em}{\MakeUppercase}
  [\vspace{1pt}\color{accent}\rule{\linewidth}{1.5pt}]
\titlespacing{\section}{0pt}{16pt}{7pt}

% ── List style ────────────────────────────────────────────────────────────────
\setlist[itemize]{
  leftmargin=1.3em, itemsep=2pt, topsep=3pt, parsep=0pt,
  label={\color{accent!60}\small\textbullet}
}
\newcommand{\tabitem}{{\color{accent!70}\textbullet}\enspace}

% ── Experience left-border box ────────────────────────────────────────────────
\newenvironment{expbox}{%
  \begin{tcolorbox}[
    blanker, breakable,
    left=10pt,
    borderline west={2.5pt}{0pt}{accentlt},
    before upper={\parskip=2pt},
    top=0pt, bottom=0pt
  ]
}{%
  \end{tcolorbox}\vspace{4pt}
}

% ── Role / Edu macros ─────────────────────────────────────────────────────────
\newcommand{\role}[4]{%
  {\bfseries #1}\hfill{\small\color{accent}#2}\\[-1pt]
  {\small\color{mid}#3}\hfill{\small\color{light}#4}\\[3pt]%
}

\newcommand{\edu}[4]{%
  {\bfseries #1}\hfill{\small\color{accent}#2}\\[-1pt]
  {\small\color{mid}#3}\hfill{\small\color{light}#4}\\[2pt]%
}

\pagestyle{empty}
\setlength{\parindent}{0pt}

% ══════════════════════════════════════════════════════════════════════════════
\begin{document}

% ── Full-bleed header ─────────────────────────────────────────────────────────
\vspace*{-1.5cm}
\begin{tcolorbox}[
  enhanced, colback=accent, colframe=accent,
  arc=0pt, boxrule=0pt,
  left=1.8cm, right=1.8cm, top=16pt, bottom=14pt,
  enlarge left by=-1.8cm, enlarge right by=-1.8cm,
  width=\textwidth+3.6cm
]
  \begin{minipage}[t]{0.54\linewidth}
    {\color{white}\LARGE\bfseries Manju Vallayil}\\[5pt]
    {\color{white!65!accent}\small
      AI Researcher\enspace·\enspace
      Explainable AI \& NLP\enspace·\enspace
      PhD, Computer Science}
  \end{minipage}%
  \hfill
  \begin{minipage}[t]{0.43\linewidth}
    \raggedleft\small\color{white!80!accent}
    \faEnvelope\enspace
      \href{mailto:manjuvallayil@gmail.com}{\textcolor{white!80!accent}{manjuvallayil@gmail.com}}\\[3pt]
    \faLinkedin\enspace
      \href{https://linkedin.com/in/manju-vallayil-phd-5b1b7a142/}{\textcolor{white!80!accent}{linkedin/manju-vallayil-phd}}\\[3pt]
    \faGlobe\enspace
      \href{https://manju.vallayil.com}{\textcolor{white!80!accent}{manju.vallayil.com}}\\[3pt]
    \faMapMarker*\enspace Auckland, New Zealand
  \end{minipage}
\end{tcolorbox}

\vspace{12pt}

% ── Profile ───────────────────────────────────────────────────────────────────
\section{Profile}

AI researcher and NLP specialist with a PhD in Computer Science, focused on Explainable AI,
context-aware evidence retrieval, and interpretable machine learning. Alongside doctoral
research, five years of independent research assistant and research officer roles
spanning AI for culturally responsive education, clinical ML for tinnitus detection, and
NLP-powered bilingual language learning technology. Each was a distinct applied engagement
with real community, healthcare, and education partners, not an extension of the thesis.
Throughout, I chose data and systems work over academic teaching positions, staying close to
applied problems rather than theoretical ones. Now seeking to bring this research depth into
industry data science and AI, where models and insights serve real organisational decisions.

% ── Key Skills ────────────────────────────────────────────────────────────────
\section{Key Skills}

\renewcommand{\arraystretch}{1.35}
\begin{tabular}{@{}p{0.48\linewidth}p{0.48\linewidth}@{}}
  \textbf{\color{mid}AI \& Machine Learning}
    & \textbf{\color{mid}Research \& Methods} \\
  \tabitem Explainable AI (XAI)
    & \tabitem NLP pipeline development \\
  \tabitem Retrieval-Augmented Generation (RAG)
    & \tabitem Literature review \& evidence synthesis \\
  \tabitem Automated fact verification
    & \tabitem Qualitative \& mixed-methods research \\
  \tabitem Thematic clustering \& embeddings
    & \tabitem Technical writing \& publication \\[20pt]
  \textbf{\color{mid}Tools \& Languages}
    & \textbf{\color{mid}Domain Knowledge} \\
  \tabitem Python (NLP, ML, data analysis)
    & \tabitem AI for education \& indigenous language \\
  \tabitem Hugging Face, scikit-learn, spaCy
    & \tabitem Human-centred \& community research \\
  \tabitem SQL, LaTeX, Git
    & \tabitem Stakeholder communication \& reporting \\
  \tabitem Microsoft Excel \& PowerPoint
    & \tabitem Research ethics \& compliance \\
\end{tabular}
\renewcommand{\arraystretch}{1}

% ── Experience ────────────────────────────────────────────────────────────────
\section{Experience}

\begin{expbox}
\role{Research Assistant — AI for Education \& Community Innovation}
     {Mar 2025 – Dec 2025}
     {Auckland University of Technology}
     {Auckland, NZ}
\begin{itemize}
  \item Evaluated AI-enabled, low-barrier tools for culturally responsive STEM education
        within the \textbf{TechTahi} project — conducting data analysis on adoption patterns
        among Māori and Pacific learners and translating findings into actionable insights.
  \item Designed and implemented research methodology to assess how AI tools are experienced
        by learners, applying mixed-methods analysis across qualitative and quantitative data.
  \item Co-authored peer-reviewed publication in \textit{Education Sciences} (May 2026),
        contributing to evidence synthesis and interpretation of AI tool outcomes.
\end{itemize}
\end{expbox}

\begin{expbox}
\role{Research Officer — AI-Driven Language Learning Technology}
     {Sep 2021 – Dec 2024}
     {Auckland University of Technology}
     {Auckland, NZ}
\begin{itemize}
  \item Co-developed \textbf{Te Kōhanga o te Tūī}, an AI-enabled bilingual learning platform
        combining speech recognition and NLP to support early childhood acquisition of
        te reo Māori and English.
  \item Contributed to applied NLP research, technical architecture documentation, and
        evaluation of speech and language model performance within a real deployment context.
  \item Liaised across language specialists, developers, and community stakeholders to
        translate research outputs into functional, accessible technology.
\end{itemize}
\end{expbox}

\begin{expbox}
\role{Research Assistant — AI Modelling for Tinnitus Detection}
     {Mar 2022 – Mar 2023}
     {Auckland University of Technology}
     {Auckland, NZ}
\begin{itemize}
  \item Applied machine learning to EEG, clinical, and psychological datasets to model
        factors influencing tinnitus detection and treatment outcomes in an interdisciplinary
        research setting.
  \item Performed data preprocessing, feature engineering, and model evaluation —
        identifying patterns across heterogeneous clinical data sources.
  \item Built and documented a minimum viable prototype (MVP) demonstrating practical
        application of AI-driven research findings to clinical stakeholders.
\end{itemize}
\end{expbox}

\begin{expbox}
\role{Business Analyst}
     {Oct 2017 – Mar 2019}
     {Tureya Limited}
     {Auckland, NZ}
\begin{itemize}
  \item Led requirements gathering and stakeholder communication for mobile application
        development — producing user needs documentation, data flow diagrams, and UX
        wireframes supporting MVP delivery.
  \item Developed and executed test plans, verifying system logic and data flows against
        specified requirements before client handover.
\end{itemize}
\end{expbox}

\begin{expbox}
\role{Software Quality Assurance Engineer}
     {May 2012 – Feb 2014}
     {TechWyse Internet Marketing}
     {Cochin, India}
\begin{itemize}
  \item Built data-driven test suites and performed SQL-based backend validation to verify
        data integrity and system accuracy across web applications.
\end{itemize}
\end{expbox}

\begin{expbox}
\role{Junior Software Quality Assurance Engineer}
     {Nov 2010 – Apr 2012}
     {ART Technology \& Software India Pvt.\ Ltd}
     {Cochin, India}
\begin{itemize}
  \item SQL-based backend data validation, test case documentation, and system testing of
        e-commerce platforms integrated with Microsoft RMS.
\end{itemize}
\end{expbox}

% ── Selected Publications ─────────────────────────────────────────────────────
\section{Selected Publications}

\begin{itemize}

  \item \textbf{Vallayil, M.}, Nand, P., Yan, W., et al.
        ``CARAG: A Context-Aware Retrieval Framework for Fact Verification, Integrating
        Local and Global Perspectives of Explainable AI.''
        \textit{Applied Sciences}, Feb 2025.
        \href{https://www.mdpi.com/2076-3417/15/4/1970}{\small[link]}

  \item \textbf{Vallayil, M.}, Nand, P., Yan, W., Allende-Cid, H.
        ``Explainability of Automated Fact Verification Systems: A Comprehensive Review.''
        \textit{Applied Sciences}, Nov 2023.
        \href{https://www.mdpi.com/2076-3417/13/23/12608}{\small[link]}

  \item \textbf{Vallayil, M.}, Nand, P., Yan, W.
        ``Explainable AI through Thematic Clustering and Contextual Visualization:
        Advancing Macro-Level Explainability in AFV Systems.''
        \textit{ACIS}, Dec 2024.
        \href{https://aisel.aisnet.org/acis2024/101/}{\small[link]}

  \item \textbf{Vallayil, M.}, Nand, P.
        \textit{manjuvallayil/factver\_master} (Hugging Face Dataset), Jan 2025.
        \href{https://doi.org/10.57967/hf/5772}{\small[link]}

\end{itemize}

{\small\color{light}Full list (10+ papers, datasets, technical reports):
\href{https://scholar.google.com/citations?user=6xGoEsMAAAAJ}{Google Scholar}\enspace·\enspace
\href{https://www.researchgate.net/profile/Manju-Vallayil-2}{ResearchGate}}

% ── Education ─────────────────────────────────────────────────────────────────
\section{Education}

\edu{Doctor of Philosophy (PhD) — Computer Science}
    {Mar 2020 – Aug 2025}
    {Auckland University of Technology (AUT)}
    {Auckland, NZ}
{\small\textit{Thesis:} \href{https://openrepository.aut.ac.nz/items/33de7437-821c-4a19-8539-0d38c6848d04}{Advancing Explainable AI: A Global Context-Aware Evidence Retrieval Framework for Interpretable Fact Verification}}

\vspace{6pt}

\edu{Master of Information Technology}
    {Mar 2016 – Apr 2019}
    {Eastern Institute of Technology}
    {New Zealand}

\edu{Bachelor of Technology (BTech) — Computer Science \& Engineering}
    {Mar 2004 – Dec 2009}
    {Government Engineering College, Palakkad}
    {India}

% ── Awards & Recognition ──────────────────────────────────────────────────────
\section{Awards \& Recognition}

\renewcommand{\arraystretch}{1.4}
\begin{tabular}{@{}p{0.78\linewidth}r@{}}
  \href{https://www.testdome.com/certificates/9fd52f85ece34b29bba43800c2918dbd}{Python — TestDome {\small(Top 10\%)}}
    & \textcolor{accent}{\small\bfseries 2026} \\
  \href{https://verify.skilljar.com/c/ngmh5b4cn9um}{Claude Code in Action Certification — Anthropic}
    & \textcolor{accent}{\small\bfseries 2026} \\
  AUT Doctoral Journal Writing Grant
    & \textcolor{accent}{\small\bfseries 2025} \\
  AUT DCT Faculty Postgraduate Summer Research Award {\small(×2)}
    & \textcolor{accent}{\small\bfseries 2023, 2024} \\
  Futures Pacific (CIIRID: IDEA) — AUT-PUCV Chile Scholarship {\small(×2)}
    & \textcolor{accent}{\small\bfseries 2021, 2023} \\
  AUT ECMS Research Scholarship
    & \textcolor{accent}{\small\bfseries 2022} \\
  PhD Full Tuition Fee Waiver
    & \textcolor{accent}{\small\bfseries 2021--2024} \\
\end{tabular}

\end{document}
